% Source: 04 -Mon Oct 5 2026.pdf, page 12. Content remains in source order.
\sourcepage{04}{12}
\section{Formula satisfiability problem (SAT)}
A boolean formula $\phi(x_1,x_2,\ldots,x_n):\bits^n\to\{0,1\}$ is a fully parenthesized expression built on
\begin{itemize}
\item $n$ variables $x_1,x_2,\ldots,x_n$;
\item binary operators: $\lor$ (OR), $\land$ (AND), $\to$ (IMPLIES), $\leftrightarrow$ (EQUALS);
\item unary operator: $\neg$ (NOT).
\end{itemize}

The new operators implies and equals : 
\begin{center}
\begin{tabular}{c|cc}
$x,y$ & $x\to y$ & $x\leftrightarrow y$\\\hline
$00$ & 1 & 1\\
$01$ & 1 & 0\\
$10$ & 0 & 0\\
$11$ & 1 & 1
\end{tabular}
\end{center}
\[
x\to y\ \equiv\ (\neg x)\lor y,\qquad
x\leftrightarrow y\ \equiv\ (x\land y)\lor(\neg x\land\neg y).
\]
\subsection{formula encoding : a tree}
We assume that the boolean formulas are always represented as their relative parenthesized trees. 
\[
\phi(x_1,x_2,x_3,x_4)=(x_1\to x_2)\lor(\neg((\neg x_1\leftrightarrow x_3)\lor x_4)).
\]
Boolean formulas can be represented by the parenthesization tree. The leafs are the  variables; the internal nodes are operators.
\sourcefigure[0.65]{assets/04/page-012-formula-tree.png}

